\section{Introduction}

Governance-style loops are awkward for off-the-shelf RL. Mean return can look fine while the policy does something brittle whenever telemetry gets adversarial or the regime flips mid-episode. At the same time, engineers rarely rip out every hand-written rule; the interesting design problem is how to bolt learning onto deterministic safeguards without pretending the safeguards do not exist.

That is the setting we had in mind for QAI-Chain. PPO still proposes an action. A cheap uncertainty read decides whether we trust that action or hand off to a deterministic heuristic. A shield---again deterministic---clips rates and bounds before anything executes. The ``quantum-inspired'' label is only about parameter count: we use a compact VQC-style readout, simulated classically, to get a nonlinear uncertainty signal without dragging a full ensemble into the loop. In the hardest regime we looked at, that readout gave higher fallback precision at the same trigger rate than MC-dropout (0.964 vs 0.743 at intensity~1.0).

We are not selling quantum speedups or a new universal RL objective. The narrow claim is that, under fast switching, calibrated routing plus a hard shield changes the failure shape relative to policy-only execution, and you can get there without a huge uncertainty module.

The empirical story is split on purpose. SafetyPointGoal1-v0 and SafetyHalfCheetahVelocity-v1 with a fixed thirty-seed protocol are the sanity check that the training stack behaves on standard constrained tasks. The governance simulators are deliberately meaner: same matched fifty-seed protocol for the controller comparison table, plus stress sweeps. Heuristic controllers stay in the plot on purpose---they are not straw opponents, and in mild regimes they still win plenty of cells.

Training changes relative to vanilla PPO are small: we add variance-aware reward shaping and a routing rule with threshold $\tau$ and risk weight $\lambda_{\mathrm{risk}}$. When we sweep complexity we keep the model class and simulator fixed and only retune those scalars so we are not confounding ``more safety glue'' with ``different backbone network.''

Blockchain commitments show up only as an optional audit trail after execution. They matter for provenance if several parties review actions; they do not change the control math we evaluate in the main experiments.

\paragraph{Contributions.}
What we think is new enough to write down:
\begin{itemize}
    \item a decision-time stack that couples uncertainty-triggered fallback with deterministic shielding,
    \item a low-footprint uncertainty head (the VQC-style block) that keeps behavior from falling apart when disturbances kick in,
    \item and an evaluation package that includes finished thirty-seed runs on the two Safety Gym tasks above, plus the fifty-seed governance controller suite and the stress tests behind the appendix figures.
\end{itemize}

\begin{table}[htbp]
\centering
\input{tables/core_contribution_layers}
\caption{Core contribution layering: what is algorithmically novel versus efficiency support versus deployment engineering.}
\label{tab:core-contribution-layers}
\end{table}

\paragraph{Scope.}
The falsifiable bit is narrow: with fixed seeds and paired comparisons, gated control plus the compact head matches the safety-relevant behavior of heavier setups at much smaller parameter budget under the non-stationary governance dynamics we actually simulated \cite{dror2018hitchhikers,pineau2021improving,wilson2017good}.

\section{Related Perspectives}

Most blockchain--ML papers we see stay in the comfortable zone: fraud scoring, mempool analytics, mechanism tuning, all with fairly stationary statistics. Adversarial governance telemetry gets less of that treatment, and reproducible benchmarks are rarer than slide decks. Nakamoto-style consensus is still the floor for why anyone trusts a shared log \cite{nakamoto2008bitcoin}; PPO is still the default hammer for continuous control \cite{schulman2017ppo}. Neither removes the need to beat on the loop before you ship it.

Safe RL has a crowded toolbox---CVaR-style objectives \cite{chow2015cvar}, robust DP views \cite{iyengar2005robust}, constrained policy optimization \cite{achiam2017cpo}, plus a long tail of distribution-shift hacks. We stayed conservative on the optimizer (PPO) and spent effort on what happens \emph{after} the policy logits are sampled: routing, shaping, shielding. Everything is evaluated under one disturbance model so the comparisons are not apples-to-oranges. We treat the paper as a systems contribution: wiring and evaluation, not a new surrogate loss.

Compared with the newer risk-sensitive and uncertainty-heavy lines, our delta is mostly architectural---how branch-level routing and a post-check shield sit in a reproducible pipeline. A wider baseline net (2023--2025 methods) is honest future work; we say that up front instead of pretending the table is exhaustive.

Reproducibility lives in the checklist-style appendix material; the main text stays focused on behavior.
